\GAGASection{4. Applications}{gaga:section:4}

\GAGASubsection{18. Algebraic character of the Betti numbers}{gaga:no:18}

Let $\sigma$ be an automorphism of the field $\mathbf{C}$. If $x$ is a point of
$\mathbf{P}_r(\mathbf{C})$ with homogeneous coordinates $t_0,\ldots,t_r$, we
write $x^\sigma$ for the point with homogeneous coordinates
$t_0^\sigma,\ldots,t_r^\sigma$. Thus $\sigma$ defines a permutation of
$\mathbf{P}_r(\mathbf{C})$.

If $X$ is a Z-closed algebraic subvariety of $\mathbf{P}_r(\mathbf{C})$, its
transform $X^\sigma$ under $\sigma$ is again a Z-closed algebraic subvariety of
$\mathbf{P}_r(\mathbf{C})$. If $X$ is non-singular, then so is $X^\sigma$ (for
example, by the Jacobian criterion).

\GAGATarget{gaga:prop:12}\textsc{Proposition 12}. --- \emph{If $X$ is
non-singular, then $X$ and $X^\sigma$ have the same Betti numbers.}

Let $b_n(X)$ be the $n$th Betti number of $X$, and let $\Omega^p(X)^h$ denote
the sheaf of germs of holomorphic differential forms of degree $p$ on $X$.
Put

\begin{equation*}
h^{p,q}(X)=\dim\mathrm{H}^q(X^h,\Omega^p(X)^h).
\end{equation*}

By the Dolbeault theorem (cf. \GAGACite{8}),

\begin{equation*}
b_n(X)=\sum_{p+q=n}h^{p,q}(X),
\end{equation*}

and similarly

\begin{equation*}
b_n(X^\sigma)=\sum_{p+q=n}h^{p,q}(X^\sigma).
\end{equation*}

By Theorem~\hyperref[gaga:thm:1]{1}, however,

\begin{equation*}
h^{p,q}(X)=\dim\mathrm{H}^q(X,\Omega^p(X)),
\end{equation*}

where $\Omega^p(X)$ now denotes the coherent algebraic sheaf of germs of
regular differential forms of degree $p$ on $X$; likewise,

\begin{equation*}
h^{p,q}(X^\sigma)=
\dim\mathrm{H}^q(X^\sigma,\Omega^p(X^\sigma)).
\end{equation*}

Moreover, if $\omega$ is a regular differential form on a Z-open subset $U$
of $X$, then $\omega^\sigma$ is regular on the Z-open subset $U^\sigma$ of
$X^\sigma$. It follows that, for every Z-open covering $\mathfrak{U}$ of $X$,
$\sigma$ defines a semilinear isomorphism from
$\mathrm{C}(\mathfrak{U},\Omega^p(X))$ onto
$\mathrm{C}(\mathfrak{U}^\sigma,\Omega^p(X^\sigma))$, hence from
$\mathrm{H}^q(\mathfrak{U},\Omega^p(X))$ onto
$\mathrm{H}^q(\mathfrak{U}^\sigma,\Omega^p(X^\sigma))$, and therefore also
from $\mathrm{H}^q(X,\Omega^p(X))$ onto
$\mathrm{H}^q(X^\sigma,\Omega^p(X^\sigma))$. Thus
$h^{p,q}(X)=h^{p,q}(X^\sigma)$, proving the proposition.

Proposition~\hyperref[gaga:prop:12]{12} gives the following result, conjectured
by A.~Weil.

\GAGATarget{gaga:prop:12:cor}\textsc{Corollary}. --- \emph{Let $V$ be a
non-singular projective variety defined over a field $K$ of algebraic numbers.
The complex varieties $X$ obtained from $V$ by embedding $K$ in $\mathbf{C}$
have Betti numbers independent of the chosen embedding.}

Indeed, any two embeddings of $K$ in $\mathbf{C}$ differ by an automorphism of
$\mathbf{C}$.

\GAGATarget{gaga:remark:18}\emph{Remark}. --- I do not know whether the
varieties $X$ and $X^\sigma$ are always homeomorphic. In any case, the example
of a curve of genus $1$ already shows that they are not always analytically
isomorphic.

\GAGASubsection{19. Chow's theorem}{gaga:no:19}

This is the following result (cf. \GAGACite{6}).

\GAGATarget{gaga:prop:13}\textsc{Proposition 13}. --- \emph{Every closed
analytic subset of projective space is algebraic.}

We show how this proposition follows from Theorem~\hyperref[gaga:thm:3]{3}.
Let $X$ be a projective space and let $Y$ be a closed analytic subset of
$X^h$. By a theorem of H.~Cartan cited above
(\hyperref[gaga:no:3]{no.~3}, Proposition~\hyperref[gaga:prop:1]{1}), the sheaf
$\mathcal{H}_Y=\mathcal{H}_X/\mathcal{A}(Y)$ is a coherent analytic sheaf on
$X^h$. Hence there is, by Theorem~\hyperref[gaga:thm:3]{3}, a coherent
algebraic sheaf $\mathcal{F}$ on $X$ such that
$\mathcal{H}_Y=\mathcal{F}^h$. By Proposition~\hyperref[gaga:prop:10:item:b]{10
b)}, the support of $\mathcal{F}^h$ is equal to that of $\mathcal{F}$ (recall,
cf. \GAGACite{18}, no.~81, that the support of $\mathcal{F}$ is the set of $x\in X$
such that $\mathcal{F}_x\ne0$), and is therefore Z-closed because
$\mathcal{F}$ is coherent. Since $\mathcal{F}^h=\mathcal{H}_Y$, this means that
$Y$ is Z-closed, as required.

We now indicate some simple applications of Chow's theorem.

\GAGATarget{gaga:prop:14}\textsc{Proposition 14}. --- \emph{If $X$ is an
algebraic variety, every compact analytic subset $X'$ of $X$ is algebraic.}

Retain the notation from the proof of Proposition~\hyperref[gaga:prop:6]{6}:
let $Y$ be a projective variety, let $U$ be a Z-open and Z-dense subset of
$Y$, and let $f:U\to X$ be a surjective regular map whose graph $T$ is
Z-closed in $X\times Y$. Put $T'=T\cap(X'\times Y)$. Since $X'$ and $Y$ are
compact and $T$ is closed, $T'$ is compact; so is its projection $Y'$ onto
the factor $Y$. On the other hand, $Y'=f^{-1}(X')$, which shows that $Y'$ is
an analytic subset of $U$, hence of $Y$. Chow's theorem therefore shows that
$Y'$ is a Z-closed subset of $Y$. Applying
Proposition~\hyperref[gaga:prop:7]{7} to $f:Y'\to X$, we conclude that
$X'=f(Y')$ is Z-closed in $X$, as required.

\GAGATarget{gaga:prop:15}\textsc{Proposition 15}. --- \emph{Every holomorphic
map $f$ from a compact algebraic variety $X$ to an algebraic variety $Y$ is
regular.}

Let $T$ be the graph of $f$ in $X\times Y$. Since $f$ is holomorphic, $T$ is
a compact analytic subset of $X\times Y$. Proposition~\hyperref[gaga:prop:14]{14}
shows that $T$ is algebraic, and hence $f$ is regular by
Proposition~\hyperref[gaga:prop:8]{8}.

\GAGATarget{gaga:prop:15:cor}\textsc{Corollary}. --- \emph{Every compact
analytic space admits at most one structure of algebraic variety.}

\GAGASubsection{20. Algebraic and analytic fibre spaces}{gaga:no:20}

Let $G$ be an algebraic group and $X$ an algebraic variety. The germs of
regular maps from $X$ to $G$ form a sheaf of groups, generally non-abelian,
which we denote by $\mathcal{G}$.

If $\mathcal{A}$ is a sheaf of groups, one can define the \emph{group}
$\mathrm{H}^0(X,\mathcal{A})$ and the \emph{set}
$\mathrm{H}^1(X,\mathcal{A})$; see \GAGACite{9} and, for example, \GAGACite{10}, Chapter~V. In
particular, $\mathrm{H}^1(X,\mathcal{G})$ is defined. Its elements are exactly
the classes of principal algebraic fibre spaces with base $X$ and structure
group $G$ (in the sense defined by A.~Weil; cf. \GAGACite{20}). For example, the
elements of $\mathrm{H}^1(X,\mathcal{O}_X)$ are the classes of fibre spaces
whose group is the additive group $\mathbf{C}$.

Likewise, if $\mathcal{G}^h$ denotes the sheaf of germs of holomorphic maps
from $X$ to $G$, the elements of
$\mathrm{H}^1(X^h,\mathcal{G}^h)$ are exactly the classes of \emph{analytic}
fibre spaces with base $X$ and group $G$. Every algebraic fibre space $E$
defines an analytic fibre space $E^h$, and therefore a map

\begin{equation*}
\varepsilon:\mathrm{H}^1(X,\mathcal{G})
\longrightarrow\mathrm{H}^1(X^h,\mathcal{G}^h),
\end{equation*}

analogous to the map in \hyperref[gaga:no:11]{no.~11}.

\GAGATarget{gaga:prop:16}\textsc{Proposition 16}. --- \emph{If $X$ is
compact, the map $\varepsilon$ is injective.}

Let $E$ and $E'$ be two principal algebraic fibre spaces with base $X$ and
structure group $G$. Proposition~\hyperref[gaga:prop:16]{16} says that, if
$E$ and $E'$ are analytically isomorphic, then they are algebraically
isomorphic. In fact, we shall prove the slightly more precise result that
every analytic isomorphism $\varphi:E\to E'$ is an algebraic isomorphism
(that is, a regular map).

The space $E\times E'$ is a principal algebraic fibre space with base
$X\times X$ and structure group $G\times G$. We denote by $(E,E')$ its
inverse image under the diagonal map $X\to X\times X$; it is the `fibred
product' of $E$ and $E'$. Let $G\times G$ act on $G$ by

\begin{equation*}
(g,g')\mathbin{.}h=gh{g'}^{-1}.
\end{equation*}

Let $T$ be the fibre space associated with the principal fibre space
$(E,E')$, with typical fibre the group $G$ equipped with the preceding
actions. It is immediate that the sections of $T$ correspond bijectively to
the isomorphisms from $E$ onto $E'$. In particular, $\varphi$ corresponds to
an analytic section $s$ of $T$. Applying Proposition~\hyperref[gaga:prop:15]{15}
to $s:X\to T$, we see that $s$ is regular. This means that $\varphi$ is
regular, and proves the proposition.

Now suppose that $X$ is a \emph{projective variety}. One may ask whether

\begin{equation*}
\varepsilon:\mathrm{H}^1(X,\mathcal{G})
\longrightarrow\mathrm{H}^1(X^h,\mathcal{G}^h)
\end{equation*}

is bijective; in other words, in view of Proposition~\hyperref[gaga:prop:16]{16},
whether every analytic fibre space is algebraic. This is plainly false if no
condition is imposed on $G$, as is shown by the case in which $G$ is an
abelian variety (or a finite group). In the following propositions we give
several groups $G$ for which it is true.

\GAGATarget{gaga:prop:17}\textsc{Proposition 17}. --- \emph{If $G$ is the
additive group $\mathbf{C}$, the map $\varepsilon$ is bijective.}

Indeed, then $\mathcal{G}=\mathcal{O}$ and
$\mathcal{G}^h=\mathcal{O}^h$, and the proposition is a special case of
Theorem~\hyperref[gaga:thm:1]{1}.

\GAGATarget{gaga:prop:18}\textsc{Proposition 18}. --- \emph{If $G$ is the
general linear group $\mathrm{GL}_n(\mathbf{C})$, the map $\varepsilon$ is
bijective.}

To every principal fibre space with structure group
$\mathrm{GL}_n(\mathbf{C})$ is associated a fibre space with vector fibre,
having typical fibre $\mathbf{C}^n$, which characterizes it. In view of the
correspondence between fibre spaces with vector fibres and locally free
sheaves (cf., for example, \GAGACite{18}, no.~41), it is therefore enough to prove the
following statement:

\emph{If $\mathcal{M}$ is a coherent analytic sheaf on $X^h$ which is locally
isomorphic to $\mathcal{H}^n$, then there is a coherent algebraic sheaf
$\mathcal{F}$ on $X$ which is locally isomorphic to $\mathcal{O}^n$ and such
that $\mathcal{F}^h$ is isomorphic to $\mathcal{M}$.}

By Theorem~\hyperref[gaga:thm:3]{3}, there is a coherent algebraic sheaf
$\mathcal{F}$ on $X$ satisfying the second condition. For every $x\in X$, the
$\mathcal{H}_x$-module

\begin{equation*}
\mathcal{F}^h_x=\mathcal{F}_x\otimes\mathcal{H}_x
\end{equation*}

is therefore isomorphic to $\mathcal{H}_x^n$. Applying \GAGAInternalRef{gaga:appendix}{Appendix}
Proposition~\GAGAInternalRef{gaga:prop:30}{30} to the rings $A=\mathcal{O}_x$, $A'=\mathcal{H}_x$, and the
module $E=\mathcal{F}_x$, we conclude that $\mathcal{F}_x$ is isomorphic to
$\mathcal{O}_x^n$. Since $\mathcal{F}$ is coherent, it follows that
$\mathcal{F}$ is locally isomorphic to $\mathcal{O}^n$, completing the proof.

\GAGATarget{gaga:prop:18:remarks}\emph{Remarks}. --- 1. For $n=1$,
$\mathrm{GL}_n(\mathbf{C})$ is the multiplicative group $\mathbf{C}^*$. If
$X$ is assumed to be a \emph{normal} variety, the group
$\mathrm{H}^1(X,\mathcal{G})$ coincides with the group of classes of divisors
locally linearly equivalent to zero (cf. \GAGACite{20}, \S~3), and
Proposition~\hyperref[gaga:prop:18]{18} says that every analytic fibre space
with base $X$ and structure group $\mathbf{C}^*$ arises from such a divisor.
When $X$ is \emph{non-singular}, this result was obtained by
Kodaira--Spencer \GAGACite{12}; in that case it is, moreover, essentially equivalent
to Lefschetz's theorem on the existence of divisors with a prescribed
homology class.

2. Proposition~\hyperref[gaga:prop:18]{18} permits other results of Kodaira to
be extended to arbitrary projective varieties, possibly with singularities;
this is so, in particular, for Theorems 7 and 8 of \GAGACite{11}. We shall not pursue
this here.

Now let $G$ be an algebraic group and $H$ an algebraic subgroup of $G$. It is
known (cf., for example, \GAGACite{13}) that the homogeneous space $G/H$ can be endowed
with the quotient algebraic-variety structure induced by that of $G$. The group $H$
acts on $G$ by right translations. We assume that these operations make $G$
into a principal algebraic fibre space with base $G/H$ and structure group
$H$, or equivalently, that there is a \emph{rational section} $G/H\to G$
(which is not always the case, as we shall see below). Under this hypothesis
we have the following result, communicated to me, together with its proof, by
A.~Grothendieck.

\GAGATarget{gaga:prop:19}\textsc{Proposition 19}. --- \emph{Let $X$ be a
compact algebraic variety, and let $P$ be a principal analytic fibre space
with base $X$ and structure group $H$. The fibre space $P$ is algebraic if and
only if the fibre space $P\times_HG$, obtained from $P$ by extending the
structure group from $H$ to $G$, is algebraic.}

Necessity is clear. To prove sufficiency, suppose that $P\times_HG$ is
algebraic. This means that there is a principal algebraic fibre space $P_0$
with structure group $G$, and an analytic isomorphism
$h:P_0\to P\times_HG$. Consider the fibre space $E$ (respectively $E_0$)
associated with $P\times_HG$ (respectively $P_0$), with typical fibre $G/H$
on which $G$ acts by translations. Then

\begin{equation*}
E_0=P_0\times_GG/H
\qquad\text{and}\qquad
E=(P\times_HG)\times_GG/H=P\times_HG/H.
\end{equation*}

The analytic isomorphism $h$ defines an analytic isomorphism $f:E_0\to E$.
But the fibre space $E=P\times_HG/H$ has a canonical section $s$, because
$H$ fixes the point of $G/H$ corresponding to the identity element of $G$.
The isomorphism $f$ transforms $s$ into a section
$s_0=f^{-1}\circ s$ of $E_0$. The section $s_0$ is holomorphic and hence,
by Proposition~\hyperref[gaga:prop:15]{15}, regular.

On the other hand, since $G$ acts on $P_0$, so does $H$, and $P_0/H$ is
precisely $E_0$. More exactly, $P_0$ is a principal algebraic fibre space with
base $E_0$ and structure group $H$. This is easily checked by a local
argument, using the hypothesis that $G$ is a principal algebraic fibre space
with base $G/H$ and structure group $H$. Let
$P_1=s_0^{-1}(P_0)$ be the inverse image of $P_0$ under the map
$s_0:X\to E_0$. The fibre space $P_1$ is a principal algebraic fibre space
with base $X$ and structure group $H$. We shall show that $P_1$ is
analytically isomorphic to $P$, which will prove the proposition.

The relation $s_0=f^{-1}\circ s$, together with the fact that $f$ is an
analytic isomorphism, shows that $P_1=s_0^{-1}(P_0)$ is analytically
isomorphic to the inverse image of $P\times_HG$, regarded as a principal
fibre space with structure group $H$, under the map $s:X\to E$. This last
inverse image is precisely $P$, as the commutative diagram shows:

\begin{center}
\begin{tikzcd}[column sep=large,row sep=large]
P \arrow[r] \arrow[d] & P\times_HG \arrow[d] \\
X \arrow[r,"s"] & E=P\times_HG/H.
\end{tikzcd}
\end{center}

This completes the proof.

Combining Propositions~\hyperref[gaga:prop:18]{18} and
\hyperref[gaga:prop:19]{19}, we obtain:

\GAGATarget{gaga:prop:20}\textsc{Proposition 20}. --- \emph{Let $G$ be an
algebraic subgroup of $\mathrm{GL}_n(\mathbf{C})$ satisfying the following
condition:}

\GAGATarget{gaga:condition:R}(R) --- \emph{There is a rational section
$\mathrm{GL}_n(\mathbf{C})/G\to\mathrm{GL}_n(\mathbf{C})$.}

\emph{Then, for every projective variety $X$, the map}

\begin{equation*}
\varepsilon:\mathrm{H}^1(X,\mathcal{G})
\longrightarrow\mathrm{H}^1(X^h,\mathcal{G}^h)
\end{equation*}

\emph{is bijective.}

\emph{Examples}. --- Condition (R) holds in the following cases:

\begin{enumerate}[label=\emph{\alph*)},leftmargin=2.2em]
\item when $G$ is \emph{solvable}, by a theorem of Rosenlicht \GAGACite{13};

\item when $G=\mathrm{SL}_n(\mathbf{C})$, the rational section then being
obvious;

\item when $G=\mathrm{Sp}_n(\mathbf{C})$, $n=2m$. In this case the homogeneous
space $\mathrm{GL}_n(\mathbf{C})/G$ is the space of non-degenerate alternating
forms

\begin{equation*}
\sum_{i<j}a_{ij}x_i\wedge x_j,
\end{equation*}

and condition (R) follows from the fact that the \emph{generic} form

\begin{equation*}
\sum_{i<j}u_{ij}x_i\wedge x_j
\end{equation*}

can be reduced to the canonical form

\begin{equation*}
\sum_{i=1}^{m}x_{2i-1}\wedge x_{2i}
\end{equation*}

by a linear change of variables with coefficients in the field
$\mathbf{C}(u_{ij})$.
\end{enumerate}

The last two examples suggest the conjecture that condition (R) holds whenever
$G$ is a \emph{simply connected semisimple group}.

On the other hand, one can show that the unimodular orthogonal group
$G=\mathrm{O}_n^+(\mathbf{C})$ does not satisfy condition (R) when $n\geqslant3$.
I do not know whether, in that case, the map

\begin{equation*}
\varepsilon:\mathrm{H}^1(X,\mathcal{G})
\longrightarrow\mathrm{H}^1(X^h,\mathcal{G}^h)
\end{equation*}

is bijective.
